%-----------------------------------LICENSE------------------------------------%
%   This file is free software: you can redistribute it and/or                 %
%   modify it under the terms of the GNU General Public License as             %
%   published by the Free Software Foundation, either version 3 of the         %
%   License, or (at your option) any later version.                            %
%                                                                              %
%   This file is distributed in the hope that it will be useful,               %
%   but WITHOUT ANY WARRANTY; without even the implied warranty of             %
%   MERCHANTABILITY or FITNESS FOR A PARTICULAR PURPOSE.  See the              %
%   GNU General Public License for more details.                               %
%                                                                              %
%   You should have received a copy of the GNU General Public License along    %
%   with this file.  If not, see <https://www.gnu.org/licenses/>.              %
%------------------------------------------------------------------------------%
%   Author: Ryan Maguire                                                       %
%   Date:   2023/12/31                                                         %
%------------------------------------------------------------------------------%
\documentclass{letter}
\signature{Ryan Maguire}
\address{%
    2-238B\\%
    Mathematics Department\\%
    MIT\\%
    Cambridge, MA USA%
}
\begin{document}
    \begin{letter}{}
        \opening{To whom it may concern,}
            I am writing to apply for the position of Astrophysicist
            (Postdoctoral Research Fellow) with the Event Horizon Telescope.
            Over the past decade, I have developed a broad and
            interdisciplinary research background in radio astronomy, Lorentz
            geometry and spacetimes, and numerical and computational
            mathematics, with a particular focus on algorithm design and
            inverse problems. I believe this combination of skills would
            make a valuable contribution to the EHT team.
            \par
            When I started working at Wellesley college with the radio science
            team, the best methods for obtaining ring profiles of the rings of
            Saturn were limited to just under 1 kilometer resolutions.
            To date, the NASA PDS still only contains 1 kilometer profiles.
            My work in this
            field has allowed us to go much further, pushing the resolution
            down to 10 to 20 meters. We are discovering new density waves in
            the rings and structures that are much finer than anything seen before.
            All of this was made possible by combining mathematics and physics
            together. We use physics to model to problem
            (using the Rayleigh-Somerfeld equation,
            which itself comes from the Maxwell equations, the wave equation, and
            the Helmholtz equation), and then we use mathematics to accurately
            solve the inverse problem to retrieve these high resolution ring
            profiles. Such a combination of skills could certainly be applied
            to analyzing data with black holes.
        \closing{Sincerely,}
    \end{letter}
\end{document}
